\section{Repository baseline and the actual formal target}

\subsection{Why another abstract calculus would be a mistake}

The central design decision is to use the repository's ordinary iPRA calculus
as the public interface.  In the pinned tree, primitive-recursive symbols are
generated by zero, successor, projections, composition, and primitive
recursion; terms are applications of those symbols to finitely scoped
variables; and \(\PRAAx\) consists of the defining equations, the successor
axioms, and induction for quantifier-free formulas.  Derivability is the
intuitionistic natural-deduction relation
\[
  \Prf\;\PRAAx\;\Gamma\;\varphi,
  \qquad\text{written in Lean as }\code{Gamma |- [PRAAx] phi}.
\]
The existing derivation library already proves substantial arithmetic in this
judgment.  The bridge transports closed and open sequents to the Foundation
presentation.  Consequently a new framework should compile into this judgment,
not ask downstream developments to adopt another proof system.

There are three different levels that must not be conflated.

\begin{enumerate}
\item A \emph{semantic metaprogram} is an ordinary Lean or Python function
which evaluates a finite object on the host natural numbers.  It is useful for
generating symbols and tests, but its correctness is not an iPRA theorem.
\item A \emph{proof-producing metacompiler} generates a finite Lean expression
whose type is a \(\Prf\;\PRAAx\) judgment.  The generator is untrusted; the
Lean kernel checks its output.
\item An \emph{object-language universal interpreter} is a primitive-recursive
program inside iPRA that consumes a code of an arbitrary formula or proof.
That is a much stronger and more cumbersome construction.  STRIVE does not
need it.
\end{enumerate}

The framework therefore specializes a fixed mathematical interface at the
metalevel.  If a theorem contains seven bounded quantifier nodes, the compiler
emits seven finite constructor applications.  The only object-level recursion
is over numerical bounds, traces, or table lengths.  This fixed-program
principle is what keeps every invariant quantifier-free.

\subsection{Checked infrastructure reused directly}

The formalization plan assumes the following existing components.

{\small\begin{longtable}{>{\raggedright\arraybackslash}p{0.34\textwidth}>{\raggedright\arraybackslash}p{0.58\textwidth}}
\toprule
Component & Role in STRIVE \\
\midrule
\endfirsthead
\toprule
Component & Role in STRIVE \\
\midrule
\endhead
\code{ConciseIPRA.PR} & Defines \code{PR}, \code{PTm}, \code{PRAAx},
instantiated defining equations, and quantifier-free induction. \\
\code{ConciseIPRA.Deduction} & Defines \code{Prf}; provides substitution,
weakening, simultaneous cut \code{Prf.trans}, and existential elimination. \\
\code{IPRADerivations.}\newline\code{Core.Basic} & Arithmetic notation, open induction helper,
instantiation of closed derivations, and congruence lemmas. \\
\code{Core.Order} and \code{Core.Semiring} & Decidable equality/order,
monotonicity, cancellation, and the finite interval algebra needed for code
bounds. \\
\code{Core.Powers} & Prefix powers in the uniform-base special case and the
algebra needed for product bounds. \\
\code{Core.MultiSums} & Primitive-recursive finite sums and products with
parameters, their defining equations, and bounded congruence. \\
\code{Core.Division} & Total quotient and remainder and uniqueness for positive
divisors. \\
\code{Core.Search} & Least bounded search with foundness and minimality
proofs. \\
\code{IPRABridge} & Closed and finite open-sequent transport to Foundation
iPRA after an ordinary derivation has been completed. \\
\bottomrule
\end{longtable}}

The use of these modules is not merely a software convenience.  It gives the
mathematical proof its intended strength.  For example, selecting the first
successful branch is not an appeal to an external choice function; it is an
application of an already formalized bounded-search recurrence.  Decoding one
digit is not an unverified use of machine division; it is reduced to the
object-theory existence and uniqueness equations for quotient and remainder.

\subsection{Status separation}

A recurring source of confusion in research packets is the phrase
``formalizable in iPRA.''  It may mean any of the following:

\begin{itemize}
\item the mathematical argument is finitary;
\item a host-language program has been tested;
\item Lean definitions have been written but contain placeholders;
\item a Lean theorem about \(\N\) has been proved;
\item an actual \(\Prf\;\PRAAx\) derivation has been checked;
\item that derivation has additionally been exported to a second presentation.
\end{itemize}

STRIVE records these statuses separately.  The present report supplies complete
human proofs and executable finite tests.  The small structural kernel is a
no-placeholder source candidate.  The arithmetic and compiler modules expose
all unfinished proof obligations as visible \code{sorry} milestones.  No new
kernel-checked theorem is asserted.

\subsection{The main unresolved engineering boundary}

Earlier framework packets already explain how to use finite queries, traces,
backward viability, stable observations, and residual producer assumptions.
Their most persistent unformalized step is a phrase of the form

\begin{quote}
At a universally controlled node, store a certificate for every response; at
an existentially controlled node, store one successful response.
\end{quote}

For a paper proof this is enough.  For \(\PRAAx\), four additional questions
must be answered.

\begin{enumerate}
\item What is the primitive-recursive bound on the whole certificate code?
\item How is a variable-length table represented by one natural number?
\item How does the checker recover the \(i\)-th subcertificate with equations
available in the object theory?
\item Which quantifier-free invariant proves that the checker has inspected
all required entries?
\end{enumerate}

The mixed-radix and signed-compiler layers below answer these questions in a
form designed around the repository's current arithmetic API.
